% GENERATED FILE. Edit canonical lessons, then run npm run generate:books.
% canonical-source-hash: fnv1a64:cf86a39c5f99a113
% canonical-lessons: PT-C442-muralha, PT-C442-fortaleza, PT-C442-mosteiro, PT-C442-claustro, PT-C442-deambular

\chapter{O centro histórico}
\label{ch:pt-o-centro-historico}
\hlchaptermodality{442}

\begin{chapteropening}
\textbf{By the end of this chapter:} \emph{I can describe an old town: the walls, the fortress, the monastery, the cloister and wandering round.}

Complete the last lesson of chapter 442: I can describe an old town: the walls, the fortress, the monastery, the cloister and wandering round.
\end{chapteropening}

\section[a muralha]{a muralha — the city wall, the rampart}

\label{lesson:PT-C442-muralha}



\begin{quote}
Before the new one: say the Portuguese for the public fountain, then the Portuguese for the pillory.
\end{quote}

\subsection*{You'll want to know: a muralha}
\textbf{a muralha} — \textquotedblleft{}the city wall, the rampart\textquotedblright{}.

\begin{quote}
A cidade velha fica dentro da muralha.
\end{quote}

\begin{quote}
\emph{The old town is inside the city wall.}
\end{quote}

An ordinary wall is \emph{o muro} or \emph{a parede}.

\begin{cousinweb}
\emph{muralha} was borrowed from Italian \emph{muraglia} 'wall', from Late Latin \emph{mūrālia} (EtymDB), built on \emph{mūrus} 'wall'.
\end{cousinweb}

\begin{grammarlens}[title={What you've built}]
One more word for telling, explaining and describing at B1.
\end{grammarlens}

\subsection*{Your turn}
\begin{itemize}
\raggedright
  \item \emph{Say it:} \emph{a muralha}
  \item \emph{Say it:} \emph{a muralha}, once more
  \item \emph{Say it:} a sentence of your own with \emph{a muralha}
  \item \emph{Recall:} say the Portuguese for the public fountain, then the Portuguese for the pillory, then say \emph{a muralha} again
\end{itemize}

\begin{tcolorbox}[breakable,colback=teal!4,colframe=teal!35!black,title={Before you move on}]
What does a muralha mean? (\textquotedblleft{}The city wall, the rampart\textquotedblright{}.) Say it once more.
\end{tcolorbox}
\section[a fortaleza]{a fortaleza — the fortress; strength}

\label{lesson:PT-C442-fortaleza}



\begin{quote}
Before the new one: say the Portuguese for the pillory, then the Portuguese for the city wall.
\end{quote}

\subsection*{You'll want to know: a fortaleza}
\textbf{a fortaleza} — \textquotedblleft{}the fortress; strength\textquotedblright{}.

\begin{quote}
A fortaleza protegia o porto contra os piratas.
\end{quote}

\begin{quote}
\emph{The fortress protected the harbour against pirates.}
\end{quote}

Also strength of mind: \emph{a fortaleza de uma pessoa}.

\begin{cousinweb}
\emph{fortaleza} came from Occitan \emph{fortalessa} (EtymDB), built on \emph{fort} 'strong', from Latin \emph{fortis}.
\end{cousinweb}

\begin{grammarlens}[title={What you've built}]
One more word for telling, explaining and describing at B1.
\end{grammarlens}

\subsection*{Your turn}
\begin{itemize}
\raggedright
  \item \emph{Say it:} \emph{a fortaleza}
  \item \emph{Say it:} \emph{a fortaleza}, once more
  \item \emph{Say it:} a sentence of your own with \emph{a fortaleza}
  \item \emph{Recall:} say the Portuguese for the pillory, then the Portuguese for the city wall, then say \emph{a fortaleza} again
\end{itemize}

\begin{tcolorbox}[breakable,colback=teal!4,colframe=teal!35!black,title={Before you move on}]
What does a fortaleza mean? (\textquotedblleft{}The fortress; strength\textquotedblright{}.) Say it once more.
\end{tcolorbox}
\section[o mosteiro]{o mosteiro — the monastery}

\label{lesson:PT-C442-mosteiro}



\begin{quote}
Before the new one: say the Portuguese for the city wall, then the Portuguese for the fortress.
\end{quote}

\subsection*{You'll want to know: o mosteiro}
\textbf{o mosteiro} — \textquotedblleft{}the monastery\textquotedblright{}.

\begin{quote}
Os monges viviam neste mosteiro.
\end{quote}

\begin{quote}
\emph{The monks lived in this monastery.}
\end{quote}

\emph{O convento} is a similar house, often for nuns.

\begin{cousinweb}
\emph{mosteiro} is inherited from Old Portuguese \emph{mõesteiro} (EtymDB), from Latin \emph{monastērium}, from Greek \emph{monast\'{\={e}}rion}, on \emph{mónos} 'alone'.
\end{cousinweb}

\begin{grammarlens}[title={What you've built}]
One more word for telling, explaining and describing at B1.
\end{grammarlens}

\subsection*{Your turn}
\begin{itemize}
\raggedright
  \item \emph{Say it:} \emph{o mosteiro}
  \item \emph{Say it:} \emph{o mosteiro}, once more
  \item \emph{Say it:} a sentence of your own with \emph{o mosteiro}
  \item \emph{Recall:} say the Portuguese for the city wall, then the Portuguese for the fortress, then say \emph{o mosteiro} again
\end{itemize}

\begin{tcolorbox}[breakable,colback=teal!4,colframe=teal!35!black,title={Before you move on}]
What does o mosteiro mean? (\textquotedblleft{}The monastery\textquotedblright{}.) Say it once more.
\end{tcolorbox}
\section[o claustro]{o claustro — the cloister}

\label{lesson:PT-C442-claustro}



\begin{quote}
Before the new one: say the Portuguese for the fortress, then the Portuguese for the monastery.
\end{quote}

\subsection*{You'll want to know: o claustro}
\textbf{o claustro} — \textquotedblleft{}the cloister\textquotedblright{}.

\begin{quote}
O claustro do mosteiro é muito silencioso.
\end{quote}

\begin{quote}
\emph{The monastery cloister is very quiet.}
\end{quote}

A square garden with covered walks round it.

\begin{cousinweb}
\emph{claustro} is a learned borrowing from Latin \emph{claustrum} 'bar, closed place' (EtymDB), from \emph{claudere} 'to close'.
\end{cousinweb}

\begin{grammarlens}[title={What you've built}]
One more word for telling, explaining and describing at B1.
\end{grammarlens}

\subsection*{Your turn}
\begin{itemize}
\raggedright
  \item \emph{Say it:} \emph{o claustro}
  \item \emph{Say it:} \emph{o claustro}, once more
  \item \emph{Say it:} a sentence of your own with \emph{o claustro}
  \item \emph{Recall:} say the Portuguese for the fortress, then the Portuguese for the monastery, then say \emph{o claustro} again
\end{itemize}

\begin{tcolorbox}[breakable,colback=teal!4,colframe=teal!35!black,title={Before you move on}]
What does o claustro mean? (\textquotedblleft{}The cloister\textquotedblright{}.) Say it once more.
\end{tcolorbox}
\section[deambular]{deambular — to wander, to stroll around}

\label{lesson:PT-C442-deambular}



\begin{quote}
Before the new one: say the Portuguese for the monastery, then the Portuguese for the cloister.
\end{quote}

\subsection*{You'll want to know: deambular}
\textbf{deambular} — \textquotedblleft{}to wander, to stroll around\textquotedblright{}.

\begin{quote}
Passámos a tarde a deambular pelas ruas antigas.
\end{quote}

\begin{quote}
\emph{We spent the afternoon wandering through the old streets.}
\end{quote}

A bookish word; it suits slow walks with no plan.

\begin{cousinweb}
\emph{deambular} is a learned borrowing from Latin \emph{deambulāre} 'to walk about' (EtymDB), built on \emph{ambulāre} 'to walk'.
\end{cousinweb}

\begin{grammarlens}[title={What you've built}]
One more word for telling, explaining and describing at B1.
\end{grammarlens}

\subsection*{Your turn}
\begin{itemize}
\raggedright
  \item \emph{Say it:} \emph{deambular}
  \item \emph{Say it:} \emph{deambular}, once more
  \item \emph{Say it:} a sentence of your own with \emph{deambular}
  \item \emph{Recall:} say the Portuguese for the monastery, then the Portuguese for the cloister, then say \emph{deambular} again
\end{itemize}

\begin{tcolorbox}[breakable,colback=teal!4,colframe=teal!35!black,title={Before you move on}]
What does deambular mean? (\textquotedblleft{}To wander, to stroll around\textquotedblright{}.) Say it once more.
\end{tcolorbox}
